\section{Introduction}
\label{sec:intro}

A call graph tells which functions a call site may invoke. Almost every
interprocedural analysis --- data-flow and taint analysis, dead-code detection,
vulnerability reachability, regression test selection, program comprehension
--- starts from one, and its errors propagate: a missing edge hides a flow, a
spurious edge creates false alarms. For dynamic languages such as Python,
constructing a static call graph is notoriously hard. Functions, classes and
bound methods are values that travel through variables, attributes, containers,
parameters and return values; methods are resolved through the C3 MRO at run
time; many calls are implicit (constructors, iterators, operators, properties,
decorators); and libraries call package code back. A symbolic analysis must
abstract all of this, and every abstraction trades precision against recall.

The state of the art is PyCG~\cite{salis2021pycg}, a whole-program,
flow-insensitive value-flow analysis that was published at ICSE 2021 and
redesigned in 2026. Large language models (LLMs) offer a different trade-off:
they read code roughly the way a developer does, use identifiers,
documentation and knowledge of popular libraries, and need no hand-written
model of the language --- but they are not sound, they hallucinate, they name
things inconsistently, and their accuracy on public code may partly be
memorisation. This raises the question that this project studies:

\begin{quote}
\emph{How does a neuro-symbolic call-graph analyzer, which lets an LLM decide
what a symbolic front-end cannot, compare with a state-of-the-art symbolic
analyzer, and what do their successes and failures reveal about both?}
\end{quote}

We answer it with a controlled comparison. As symbolic baselines we run both
versions of PyCG unmodified (Section~\ref{sec:prelim}). As the neuro-symbolic
solution we design and implement \nscg (Section~\ref{sec:approach}): a symbolic
front-end names the nodes exactly as the ground truth does, enumerates 13 kinds
of call sites including implicit ones, pre-resolves callees with a cheap
resolver and proposes candidates by name, callbacks, argument flows and
library hooks; one LLM request per chunk of sites decides the callees; and a
symbolic validator checks the answers against the definitions, maps classes to
constructors through the MRO, repairs wrong paths and rejects hallucinated
names. Several variants isolate the contribution of each component, including
an LLM-only analyzer (V0) and a hybrid that gives PyCG's answer to the LLM (V3).
We evaluate all tools with the official scorer of the PyCG evaluation on its
micro-benchmark (622 programs, one per language feature) and on four real
packages whose ground truths were written by hand and audited four times
(Section~\ref{sec:setup}).\looseness=-1

\textbf{Findings.} The comparison yields five main results:
\begin{enumerate}[leftmargin=*,nosep]
\item \textbf{Parity on real packages.} On the four packages, \nscg
(\R{macro-ins-f1-nscg-v21} $F_1$ inside the package) is on par with PyCG-2026
(\R{macro-ins-f1-pycg-2026}) and far ahead of PyCG-2021
(\R{macro-ins-f1-pycg-2021}); the tools make different false positives, so
their intersection is the most precise and their union the most complete
configuration on the packages.
\item \textbf{More complete, less exact.} On the micro-benchmark, the LLM-based
analyzers have the highest recall but lower precision and about half as many
exactly correct programs as PyCG-2026, because they over-approximate value flow
and misname outside callees (Section~\ref{sec:rq1}).\looseness=-1
\item \textbf{The LLM alone is strong but brittle.} V0 is a strong baseline, and
a renaming probe (Section~\ref{sec:rq4}) shows that its micro accuracy does not
depend on the original identifiers; but V0 cannot be checked, and with a
stronger model it gets \emph{worse}, while \nscg gets better (Section~\ref{sec:rq2}).\looseness=-1
\item \textbf{The LLM filters, it does not discover.} Nearly all of \nscg's
missed package edges were never proposed by the front-end, and most of its
spurious ones were proposed by name and wrongly accepted by the LLM (Section~\ref{sec:rq3}).\looseness=-1
\item \textbf{A shared hard core.} All tools miss the same edges: callbacks
from library code (Fabric on top of Invoke) and functions that the
documentation assigns to a parameter's role (funcy).
\end{enumerate}

\textbf{Contributions.} (i) \nscg, an original neuro-symbolic call-graph
analyzer with a symbolic front-end and validator around a constrained LLM step;
(ii) a fair comparison with two versions of PyCG on audited real-world ground
truths, with ablations of every component, repetitions, a second model, and a
memorisation probe; (iii) an error taxonomy that explains the strengths and
limits of both approaches; and (iv) a reproducible artifact in which every LLM
answer is cached, available at \repo.
